\documentclass{article}
\usepackage[T1]{fontenc}
\usepackage{lmodern}
\usepackage{graphicx}
\usepackage{amsmath,amssymb}
\usepackage[a4paper,margin=2cm]{geometry}
\usepackage[absolute,overlay]{textpos}
\setlength{\TPHorizModule}{1mm}
\setlength{\TPVertModule}{1mm}
\usepackage[indonesian]{babel}
\usepackage{hyperref}
\setlength{\emergencystretch}{2em}
% unit: Halaman 11

\begin{document}

% === Halaman 11 (llm) ===

Karena bobot total tersisa = 0, maka solusi persoalan \textit{superincreasing knapsack} ditemukan. Barisan bobot yang dimasukkan ke dalam \textit{knapsack} adalah

\[ \{2, 3, -, 13, -, 52\} \]

sehingga

\[ 70 = (1 \times 2) + (1 \times 3) + (0 \times 6) + (1 \times 13) + (0 \times 27) + (1 \times 52) \]

Dengan kata lain, plainteks dari kriptogram 70 adalah \textcolor{red}{110101}.

\end{document}
